\subsection{Latency Test Workflow Pseudocode}\label{sec:latency_workflow_pseudocode}

\vspace{0.20cm}

\RestyleAlgo{ruled}
\LinesNumbered
\SetAlgoLined
\DontPrintSemicolon
\SetAlgoSkip{medskip}
\SetAlCapSkip{0.50em}
\SetAlgoInsideSkip{smallskip}
\SetAlgoCaptionSeparator{\quad}

\subsubsection{System Latency Experiments}\label{sec:latency_system_pseudocode}

\vspace{0.30cm}

\begin{algorithm}[H]
\small
\setlength{\baselineskip}{1.15\baselineskip}
\caption{Arduino direct-servo end-to-end latency test}
\label{alg:latency_arduino_e2e}
\KwIn{Uploaded direct-servo logger firmware, command profile, logger configuration, logic-analyser configuration, and optional existing logger folder.}
\KwOut{Matched end-to-end latency events and per-surface, overall, and integrity summaries.}
Prepare the Nano 33 IoT logger firmware in Arduino IDE and verify the board is reachable on the logger transport\;
\eIf{a hardware run is requested}{
    Configure the repeatable step-train command profile and logic-analyser capture\;
    Start the analyser capture, run the Arduino logger test, and store the host and board logs\;
}{
    Select the requested or latest logger folder for import-only analysis\;
}
Load host dispatch records, board receive and apply records, profile events, and analyser reference and servo-output captures\;
Map the host, board, and analyser clocks onto a common time base\;
Match each commanded transition to a valid output pulse transition using the configured association windows\;
\For{each matched or candidate event}{
    Reject events with missing records, ambiguous transitions, invalid ordering, or failed physical consistency checks\;
    Compute the latency segments from command scheduling through board application and servo-output response\;
}
Summarise valid events by control surface and over the full run, and record integrity counts for rejected and unmatched events\;
\textbf{return} the event table and latency audit summaries\;
\end{algorithm}

\vspace{0.70cm}

\begin{algorithm}[H]
\small
\setlength{\baselineskip}{1.15\baselineskip}
\caption{Nano transmitter end-to-end latency test}
\label{alg:latency_transmitter_nano_e2e}
\KwIn{Uploaded Nano transmitter firmware, transmitter command profile, host logger configuration, capture inputs, clock-repair settings, and Python post-processor settings.}
\KwOut{Trainer and receiver transition events, clock-repair audit, and transmitter-chain latency summaries.}
Prepare the Nano 33 IoT transmitter firmware in Arduino IDE and verify the serial and trainer-output paths\;
\eIf{a hardware run is requested}{
    Configure the repeatable transmitter step train, run the host-to-Nano transmitter test, and capture reference, trainer, and receiver signals\;
}{
    Select the requested or latest transmitter logger folder for import-only analysis\;
}
Load host dispatch logs, board receive and commit logs, captured pulse trains, and recorded profile events\;
Repair the Nano microsecond clock offline so board events can be compared with host and capture times\;
Run the post-processor when outputs are missing, forced, or stale\;
Detect stable command states and output transitions in the trainer and receiver captures\;
\For{each commanded transition}{
    Associate scheduled, host-dispatch, board-commit, trainer-output, and receiver-output events where the transition is unique\;
    Reject missing, duplicated, or timing-inconsistent matches\;
}
Summarise the retained latency chain from scheduling and dispatch through transmitter commit, trainer output, and receiver response\;
\textbf{return} the matched event table, clock audit, and latency summaries\;
\end{algorithm}

\cleardoublepage

\phantom{.}

\vspace{0.50cm}

\subsubsection{Overall Latency Experiments}\label{sec:latency_overall_pseudocode}

\vspace{0.30cm}

\begin{algorithm}[H]
\small
\setlength{\baselineskip}{1.15\baselineskip}
\caption{Overall Vicon-corrected bang-bang latency test}
\label{alg:latency_overall_bang_bang}
\KwIn{Overall latency configuration, Nano serial link, Vicon tracker setup, bang-bang command profile, selected Vicon filter, and post-analysis thresholds.}
\KwOut{Raw control-path logs, accepted latency events, rejected-event audit, and per-surface latency summaries.}
Configure the bang-bang command schedule, neutral intervals, Vicon state filter, and run label\;
Open the Nano serial link, command neutral, connect to Vicon, and calibrate the neutral state\;
\For{each control period until the scheduled run ends}{
    Read the latest Vicon rigid-body and surface state samples\;
    Drain serial telemetry from the Nano control path\;
    \If{a command update is due}{
        Evaluate the bang-bang profile, encode the command packet, write it to the Nano, and log schedule and write timestamps\;
    }
}
Return the surfaces to neutral, drain remaining telemetry, close the hardware links, and save the raw run logs\;
Extract the event table, command log, filtered Vicon state log, and serial telemetry for post-analysis\;
\For{each commanded bang-bang event}{
    Estimate the pre-command baseline and final response levels from Vicon windows\;
    Find the response crossing times at the selected motion fractions\;
    Reject events with missing windows, poor tracking quality, small response amplitude, or inconsistent crossings\;
}
Summarise accepted and rejected events by surface and over the full test\;
\textbf{return} the raw logs, latency event table, rejected-event table, and summary audit\;
\end{algorithm}

\vspace{0.70cm}

\begin{algorithm}[H]
\small
\setlength{\baselineskip}{1.15\baselineskip}
\caption{Nausicaa rigid-body Vicon motion characterization}
\label{alg:latency_vicon_motion}
\KwIn{Vicon subject name, static and dynamic phase schedule, sample period, tracker configuration, and candidate filter set.}
\KwOut{Raw and processed pose logs, phase statistics, filter-candidate table, and recommended filter setting.}
Configure Vicon to record the target rigid body without requiring control-surface subjects\;
Connect to the tracker and collect pose samples through the scheduled static and dynamic phases\;
Close the tracker and reject the run if no valid subject samples are available\;
Sort the raw samples by time, assign phase labels, and remove occluded or unusable rows from the processed trajectory\;
Compute position, attitude, velocity, angular-rate, sample-quality, and high-frequency residual statistics for each phase\;
\For{each candidate filter}{
    Apply the filter to the processed trajectory and estimate delay, residual noise, and dynamic-motion preservation\;
    Score the candidate using the static-noise and dynamic-delay trade-off required by the latency tests\;
}
Select the lowest-scoring acceptable filter and export the phase and filter audits\;
\textbf{return} the motion-characterization logs and recommended filter setting\;
\end{algorithm}
